\section{Post-Estimation}
\label{sec:post_estimation}

Every estimator in EconIRL returns the same \texttt{EstimationSummary} object, which feeds into a shared infrastructure for inference, diagnostics, goodness-of-fit evaluation, and counterfactual simulation.
This section describes each component.

\paragraph{Standard errors.}
EconIRL provides four methods for computing standard errors, selected via the \texttt{se\_method} parameter on \texttt{.estimate()}.
Asymptotic standard errors use the inverse of the negative Hessian, $\mathrm{Var}(\hat\theta) = (-H)^{-1}$, which is valid under correct specification of the likelihood.
Robust sandwich standard errors compute $H^{-1}BH^{-1}$ where $B = \sum_i g_i g_i^\top$ is the outer product of per-observation score vectors, accounting for potential misspecification without requiring a correctly specified variance.
Bootstrap standard errors resample individuals with replacement via \texttt{TrajectoryPanel.resample\_individuals()} and re-estimate on each sample.
The warm-start one-step NPL bootstrap runs only one to three Newton steps from the MLE estimate on each bootstrap sample, achieving quadratic convergence with roughly 100x speedup over full re-estimation.
Clustered standard errors aggregate the score vectors within individuals before forming the outer product, with a small-sample correction factor of $G/(G-1) \times (N-1)/(N-K)$ where $G$ is the number of clusters, $N$ is the number of observations, and $K$ is the number of parameters.
The Hessian is computed analytically for NFXP via implicit differentiation through the Bellman fixed point \citep{iskhakov2016} and numerically via central finite differences for all other estimators.

\paragraph{Identification diagnostics.}
Post-estimation, the library performs eigenvalue analysis of the negative Hessian at the optimum.
The \texttt{check\_identification()} function classifies the result into one of four categories based on eigenvalue thresholds.
If the rank of $-H$ is less than $K$, the model is under-identified.
If any eigenvalue is non-positive, the optimum is a saddle point or local minimum.
If the condition number exceeds $10^6$, the model is weakly identified.
Otherwise the model is well-identified.
Parameter correlations exceeding 0.95 are flagged as a warning of near-collinearity \citep{magnacThesmar2002}.
Parameter sensitivity analysis computes the elasticity of each parameter estimate with respect to all other parameters, revealing which parameters are tightly coupled and which are independently identified by the data.

\paragraph{Goodness of fit.}
The \texttt{GoodnessOfFit} object reports log-likelihood, AIC, BIC, McFadden pseudo-$R^2$, and prediction accuracy, defined as the fraction of observations where the modal predicted action matches the observed action.
Four additional fit metrics are available.
The Brier score measures calibration of the predicted choice probabilities, with values ranging from 0 (perfect) to 2 (worst possible).
The KL divergence between data and model CCPs, weighted by empirical state frequencies, measures the information loss from using the model distribution in place of the empirical distribution.
Efron pseudo-$R^2$ provides a variance-ratio analog of the linear $R^2$ for discrete choice settings.
A Pearson $\chi^2$ CCP consistency test compares the observed and predicted choice frequencies at each state, returning a test statistic, degrees of freedom, and $p$-value for the null hypothesis that the model generates the data.

\paragraph{EstimationSummary.}
The \texttt{EstimationSummary} object provides parameter estimates, standard errors, $t$-statistics, $p$-values, and 95 percent confidence intervals.
Joint linear hypotheses $H_0\!: R\theta = r$ are tested via Wald statistics with associated degrees of freedom and $p$-values.
The \texttt{.summary()} method prints a formatted table in the style of the \texttt{statsmodels} library, with one row per parameter.
The \texttt{.to\_latex()} method produces a publication-ready \texttt{booktabs} table.
The \texttt{.to\_dataframe()} method returns a pandas DataFrame for programmatic analysis.
This uniform output format means that switching from NFXP to MCE-IRL or from CCP to AIRL requires changing only the estimator class, with no changes to downstream inference or reporting code.

\paragraph{Counterfactual simulation.}
The library supports four types of counterfactual analysis, following the taxonomy of \citet{aguirregabiriaMira2010}.
Type 1 (state extrapolation) evaluates the estimated policy at shifted state indices without re-solving the Bellman equation, answering questions about what happens when the initial state distribution changes.
Type 2 (environment change) re-solves the Bellman equation under modified transition dynamics while holding the reward function fixed, answering questions like ``what if component reliability improves by 20 percent.''
Type 3 (reward change) re-solves under modified utility parameters while holding transitions fixed, answering ``what if replacement cost doubles.''
Type 4 (welfare decomposition) decomposes the total welfare change into contributions from each channel.
Elasticity analysis sweeps a single parameter over a user-specified range, re-solves for the optimal policy at each point, and reports the policy response surface.
The \texttt{simulate\_counterfactual()} function generates synthetic panels under both baseline and counterfactual policies, producing action frequency and state frequency comparisons for direct visualization.
Stationary distributions under each policy are computed via power iteration to assess long-run behavioral effects.
